\documentclass[11pt]{article}
\usepackage[preprint]{acl}
\usepackage{times}
\usepackage{latexsym}
\usepackage[T1]{fontenc}
\usepackage[utf8]{inputenc}
\usepackage{microtype}
\usepackage{inconsolata}
\usepackage{graphicx}
\usepackage{amsmath}

% Slightly enlarge the title box to fit authors cleanly.
\setlength\titlebox{5.0cm}

\title{Retrieval-Based Scientific Claim Verification with SciFact}

\author{Jiapei Liao \\
  Boston University \\
  \texttt{jiapei03@bu.edu}
  \And
  Yuxiang Liu \\
  Boston University \\
  \texttt{mat1900@bu.edu}}

\begin{document}
\maketitle

\begin{abstract}
This proposal describes an \textbf{empirical project} on scientific claim verification. Given a scientific claim and a closed corpus of biomedical paper abstracts, we aim to retrieve evidence-bearing abstracts and predict whether the retrieved evidence \textsc{Supports}, \textsc{Refutes}, or provides \textsc{NoInfo} for the claim. We will use SciFact, which contains 1,409 claims and 5,183 biomedical abstracts, as a manageable but realistic benchmark for evidence-grounded scientific reasoning \citep{wadden2020fact}. Our plan is to compare a simple lexical pipeline against a stronger dense-retrieval-plus-encoder pipeline, using sentence embeddings for retrieval and SciBERT or RoBERTa for claim verification. We will evaluate both component-level and end-to-end performance to determine whether gains come mainly from better retrieval or better stance prediction.
\end{abstract}

\section{Introduction}
This proposal outlines an \textbf{empirical project} on scientific claim verification. We will study how to build a system that reads a scientific claim, searches a corpus of paper abstracts, and predicts whether the retrieved evidence \textsc{Supports} the claim, \textsc{Refutes} it, or provides \textsc{NoInfo}. We will use the SciFact dataset because it offers a realistic benchmark for evidence-based scientific reasoning while still being manageable for a course project. SciFact is a closed-corpus benchmark with 1,409 expert-written scientific claims, a corpus of 5,183 biomedical abstracts, and a standard split of 809 training claims, 300 development claims, and 300 test claims \citep{wadden2020fact}. To keep the project scope realistic, we will focus on abstract retrieval and abstract-level stance classification; sentence-level rationale extraction will be treated as a possible extension rather than a required first-stage goal.

\subsection{What Are We Trying To Do?}
Our project studies a central NLP question: can a model retrieve evidence that is not only topically related to a claim, but actually justifies or refutes it? This task is important because scientific search is not the same as keyword search. An abstract may mention the same terms as a claim and still fail to provide useful evidence, while another abstract with different wording may directly support or refute it. SciFact was introduced specifically to study this setting of evidence-grounded scientific claim verification \citep{wadden2020fact}.

Formally, let $c$ be a claim and let $\mathcal{D}$ be a corpus of scientific abstracts. Our system has two stages. First, a retriever returns the top-$k$ candidate abstracts,
\begin{equation}
R_k(c) = \operatorname{TopK}_{a \in \mathcal{D}}\; \mathrm{sim}(g(c), g(a)),
\end{equation}
where $g(\cdot)$ maps text into an embedding space. Then, for each retrieved pair $(c,a)$, a classifier predicts a label,
\begin{equation}
\hat{y} = f_{\theta}(c,a), \qquad \hat{y} \in \{\textsc{Supports},\textsc{Refutes},\textsc{NoInfo}\}.
\end{equation}
The system succeeds only if it retrieves useful evidence and assigns the correct stance label.

\subsection{Example of Input and Output}
Our project has two clear components.

\paragraph{1. The Evidence Retriever}
This module searches the SciFact corpus for abstracts that may contain evidence.
\begin{itemize}
    \item \textbf{Input:} A natural-language scientific claim and the full abstract corpus.
    \item \textbf{Output:} A ranked top-$k$ list of candidate abstracts.
\end{itemize}

\paragraph{2. The Claim Verifier}
This module decides whether a retrieved abstract supports the claim, refutes it, or contains no real evidence.
\begin{itemize}
    \item \textbf{Input:} A claim--abstract pair $(c,a)$.
    \item \textbf{Output:} One label from \textsc{Supports}, \textsc{Refutes}, or \textsc{NoInfo}.
\end{itemize}

For example, if the claim is ``Drug X reduces the risk of disease Y,'' the retriever should find relevant abstracts about Drug X and disease Y, and the verifier should judge whether each abstract actually supports, refutes, or fails to verify that statement.

\subsection{How Will We Build and Evaluate the System?}
We will use the SciFact corpus and claim splits as our dataset \citep{wadden2020fact}. Our software stack will include Python, PyTorch, Hugging Face Transformers, Sentence-Transformers, FAISS, and scikit-learn. For dense retrieval, we plan to use a sentence embedding model such as Sentence-BERT with cosine similarity over a FAISS index \citep{reimers2019sentence,johnson2017billion}. For verification, we plan to fine-tune a pretrained encoder such as SciBERT or RoBERTa on claim--abstract pairs \citep{beltagy2019scibert,liu2019roberta}. If full fine-tuning is too expensive, we will freeze the encoder and train a lightweight classifier on top.

Our \textbf{baseline} will be simple and easy to implement: the official SciFact TF--IDF retriever, followed by a multinomial logistic regression classifier over TF--IDF features of the concatenated claim and abstract. This gives us a transparent lexical baseline for the full pipeline.

Our main system will replace both stages with stronger pretrained models. For retrieval, we plan to use a sentence embedding model with cosine similarity over a FAISS index for dense retrieval. For verification, we plan to fine-tune a pretrained encoder such as SciBERT or RoBERTa on claim--abstract pairs. We will compare whether dense retrieval improves over TF--IDF and whether a scientific-domain or strong general-domain encoder improves over the lexical baseline.

We will evaluate retrieval with Recall@$k$ (for example, $k \in \{3,5,10\}$), mean reciprocal rank, and evidence hit rate. We will evaluate classification with macro-F1, per-class F1, and accuracy. Because this is a pipeline, we will also report an end-to-end metric: whether a gold evidence abstract appears in the top-$k$ retrieved set and receives the correct label. We will consider the project successful if dense retrieval improves evidence recall over TF--IDF and the encoder-based verifier improves macro-F1 over the lexical baseline.

\subsection{How Will This Project Be Helpful?}
This project is useful for three reasons:
\begin{itemize}
    \item \textbf{Improving Scientific Search:} It studies how NLP systems can move from topical retrieval to evidence-based claim verification.
    \item \textbf{Understanding Errors:} Because the system is modular, we can separate retrieval failures from classification failures and analyze them clearly.
    \item \textbf{Building a Feasible End-to-End Project:} SciFact is realistic enough to be meaningful, but compact enough for controlled experimentation in a course setting.
\end{itemize}

\section{Description of Progress}
\subsection{Methods}
We have implemented the two-stage pipeline proposed above and evaluated both component-level and end-to-end performance. Our current baseline retriever is a TF--IDF system over the concatenated title and abstract of each SciFact document. Our stronger retriever is a dense semantic retriever based on \texttt{sentence-transformers/all-mpnet-base-v2}, which encodes claims and abstracts into a shared embedding space and ranks documents by vector similarity. In this midway stage, the dense retriever is used as a fixed pretrained encoder rather than being fine-tuned on SciFact \citep{song2020mpnet}.

For verification, we implemented two document-level classifiers. The baseline verifier is multinomial logistic regression over TF--IDF features extracted from the concatenated claim and abstract. Our stronger verifier is \textsc{SciBERT} fine-tuned for three-way claim verification (\textsc{Supports}, \textsc{Refutes}, \textsc{NoInfo}) \citep{beltagy2019scibert}. We fine-tune \texttt{allenai/scibert\_scivocab\_uncased} on claim--abstract pairs using class-weighted cross-entropy, mixed-precision GPU training, and checkpoint selection by development macro-F1. The final model was trained for up to 10 epochs, and the best checkpoint appeared at epoch 6.

To train the verifier, we construct document-level pairs from the SciFact annotations. Every gold evidence document for a claim becomes a positive pair labeled \textsc{Supports} or \textsc{Refutes}. We then add four \textsc{NoInfo} negatives per claim, first from cited-but-non-evidence documents, then from TF--IDF hard negatives, and finally from random corpus negatives when needed. This gives the verifier a realistic mixture of easy and difficult negative examples.

\subsection{Experiments}
We use the official SciFact train and development splits \citep{wadden2020fact}. The corpus contains 5,183 abstracts. The train split contains 809 claims, and the development split contains 300 claims. After pair construction, the verifier is trained on 3,800 claim--document pairs and evaluated on 1,409 development pairs. The label distribution is imbalanced, with \textsc{NoInfo} as the majority class, so macro-F1 is especially important for model comparison.

We report three kinds of evaluation. For retrieval, we measure Recall@$k$ for $k \in \{3,5,10\}$, mean reciprocal rank (MRR), and evidence coverage. For oracle document classification, we measure accuracy, macro-F1, and per-class F1. For full-pipeline evaluation, we report \textit{JointHit@$k$}, which counts a claim as correct if at least one gold evidence document appears in the top-$k$ retrieved set and also receives the correct stance label. This lets us distinguish whether an improvement comes from retrieval, verification, or both.

All reported midway experiments were run on a single GPU with mixed precision. The stronger verifier used a train batch size of 16, an evaluation batch size of 32, a learning rate of $2\times10^{-5}$, and a maximum sequence length of 256. Training the \textsc{SciBERT} verifier took about 412 seconds in total.

\begin{figure}[t]
\centering
\begin{minipage}{0.48\linewidth}
\centering
\includegraphics[width=\linewidth]{midway_figures/retrieval_recall_curve.png}
\end{minipage}
\hfill
\begin{minipage}{0.48\linewidth}
\centering
\includegraphics[width=\linewidth]{midway_figures/transformer_confusion_matrix.png}
\end{minipage}
\caption{Left: development-set retrieval recall for TF--IDF and dense retrieval. Right: confusion matrix for the fine-tuned \textsc{SciBERT} verifier on oracle document classification.}
\label{fig:results}
\end{figure}

\begin{table}[t]
\centering
\small
\begin{tabular}{lcccc}
\hline
\textbf{Retriever} & \textbf{R@3} & \textbf{R@5} & \textbf{R@10} & \textbf{MRR} \\
\hline
TF--IDF & 0.761 & 0.803 & 0.856 & 0.657 \\
Dense (\texttt{all-mpnet-base-v2}) & \textbf{0.798} & \textbf{0.872} & \textbf{0.931} & \textbf{0.743} \\
\hline
\end{tabular}
\caption{Retrieval performance on the SciFact development set.}
\label{tab:retrieval}
\end{table}

\begin{table}[t]
\centering
\small
\resizebox{\linewidth}{!}{
\begin{tabular}{lcccc}
\hline
\textbf{Verifier / Pipeline} & \textbf{Accuracy} & \textbf{Macro-F1} & \textbf{JointHit@3} & \textbf{JointHit@5} \\
\hline
TF--IDF + Logistic Regression & 0.836 & 0.359 & 0.048 & 0.053 \\
TF--IDF + \textsc{SciBERT} & --- & --- & 0.452 & 0.479 \\
Dense + Logistic Regression & --- & --- & 0.053 & 0.059 \\
Dense + \textsc{SciBERT} & --- & --- & \textbf{0.484} & \textbf{0.521} \\
\hline
\textsc{SciBERT} oracle verifier & \textbf{0.913} & \textbf{0.676} & --- & --- \\
\hline
\end{tabular}
}
\caption{Oracle verification and end-to-end pipeline results. The TF--IDF + Logistic Regression line serves as the full lexical baseline.}
\label{tab:pipeline}
\end{table}

\subsection{Results}
Our main midway result is that the stronger models improve both retrieval and verification, and these gains transfer to the full pipeline. As shown in Table~\ref{tab:retrieval}, dense retrieval outperforms TF--IDF at every cutoff, with especially clear gains at Recall@5, Recall@10, and MRR. This suggests that semantic matching is more effective than pure lexical overlap for finding evidence-bearing abstracts in SciFact.

The verifier comparison shows an even larger gap. The lexical verifier achieves reasonable overall accuracy because it handles the dominant \textsc{NoInfo} class fairly well, but its macro-F1 is low and it performs poorly on \textsc{Supports} and especially \textsc{Refutes}. In contrast, the fine-tuned \textsc{SciBERT} verifier substantially improves macro-F1 and produces a much more balanced confusion matrix (Figure~\ref{fig:results}). This is encouraging because it suggests that domain-specific pretrained encoders are already useful at the abstract level, even without rationale extraction.

The end-to-end results in Table~\ref{tab:pipeline} confirm that both stages matter. Replacing only the retriever while keeping the lexical verifier gives almost no gain, which indicates that verification is the main bottleneck in the baseline pipeline. Replacing the verifier with \textsc{SciBERT} gives a large jump, and combining dense retrieval with \textsc{SciBERT} gives the strongest overall system. The best system reaches JointHit@5 of 0.521, compared with 0.053 for the full lexical baseline. A remaining weakness is that \textsc{Refutes} examples are still harder than \textsc{Supports}, and development performance peaks around epoch 6 before fluctuating, which suggests mild overfitting on this small dataset.

\section{Project Plan}
For the final report, we plan to extend this midway system in three directions. First, we will perform more detailed error analysis to separate retrieval misses from stance classification mistakes and to understand which claims are most difficult for the current verifier. Second, we will explore one stronger extension beyond abstract-level verification, most likely sentence-level rationale extraction or a more targeted evidence selection stage, since the current system still treats entire abstracts as evidence units. Third, we will run one additional model comparison, likely \textsc{PubMedBERT} or a stronger biomedical verifier, to test whether the remaining gap comes from domain mismatch or from the lack of rationale supervision.

Compared with the original proposal, our trajectory has become more concrete. We have already implemented the lexical baseline and a stronger dense-retrieval-plus-\textsc{SciBERT} pipeline, and we now know that verification is the larger source of improvement. The final stage of the project will therefore focus less on adding many unrelated baselines and more on strengthening the evidence modeling step and explaining where the current best system still fails.

\begin{thebibliography}{9}

\bibitem[Beltagy et~al.(2019)Beltagy, Lo, and Cohan]{beltagy2019scibert}
Iz Beltagy, Kyle Lo, and Arman Cohan. 2019.
\newblock {SciBERT}: A pretrained language model for scientific text.
\newblock In \textit{Proceedings of the 2019 Conference on Empirical Methods in Natural Language Processing and the 9th International Joint Conference on Natural Language Processing (EMNLP-IJCNLP)}, pages 3615--3620.

\bibitem[Johnson et~al.(2017)Johnson, Douze, and J\'egou]{johnson2017billion}
Jeff Johnson, Matthijs Douze, and Herv\'{e} J\'{e}gou. 2017.
\newblock Billion-scale similarity search with GPUs.
\newblock \textit{arXiv preprint arXiv:1702.08734}.

\bibitem[Liu et~al.(2019)Liu, Ott, Goyal, Du, Joshi, Chen, Levy, Lewis, Zettlemoyer, and Stoyanov]{liu2019roberta}
Yinhan Liu, Myle Ott, Naman Goyal, Jingfei Du, Mandar Joshi, Danqi Chen, Omer Levy, Mike Lewis, Luke Zettlemoyer, and Veselin Stoyanov. 2019.
\newblock RoBERTa: A robustly optimized {BERT} pretraining approach.
\newblock \textit{arXiv preprint arXiv:1907.11692}.

\bibitem[Reimers and Gurevych(2019)]{reimers2019sentence}
Nils Reimers and Iryna Gurevych. 2019.
\newblock Sentence-BERT: Sentence embeddings using siamese {BERT}-networks.
\newblock In \textit{Proceedings of the 2019 Conference on Empirical Methods in Natural Language Processing and the 9th International Joint Conference on Natural Language Processing (EMNLP-IJCNLP)}, pages 3982--3992.

\bibitem[Song et~al.(2020)Song, Tan, Qin, Lu, and Liu]{song2020mpnet}
Kaitao Song, Xu Tan, Tao Qin, Jianfeng Lu, and Tie-Yan Liu. 2020.
\newblock MPNet: Masked and permuted pre-training for language understanding.
\newblock In \textit{Advances in Neural Information Processing Systems}, pages 16857--16867.

\bibitem[Wadden et~al.(2020)Wadden, Lin, Lo, Wang, van Zuylen, Cohan, and Hajishirzi]{wadden2020fact}
David Wadden, Shanchuan Lin, Kyle Lo, Lucy Lu Wang, Madeleine van Zuylen, Arman Cohan, and Hannaneh Hajishirzi. 2020.
\newblock Fact or fiction: Verifying scientific claims.
\newblock In \textit{Proceedings of the 2020 Conference on Empirical Methods in Natural Language Processing (EMNLP)}, pages 7534--7550.

\end{thebibliography}

\end{document}
